\documentclass[]{article}
\usepackage{amsmath}      % For advanced math, gather, align, and cases
\usepackage{amssymb}      % For mathematical symbols
\usepackage{graphicx}     % For including images
\usepackage{subcaption}   % For subfigures (Layout 6.1)
\usepackage{multirow}     % For multirow table cells
\usepackage{makecell}     % For multi-line text inside a single table cell
\usepackage{physics} % \dv, \pdv, \ket (optional)
\usepackage{float}

\graphicspath{{./}}
\title{Mathematical Foundations of Image Synthesis}
\author{2305009}
\date{January 4, 2026}
\begin{document}
\maketitle
\section{Introduction}
Image synthesis is the process of generating images from mathematical and algo-
rithmic descriptions. Modern graphics systems rely on 
\textbf{precise mathematical models}
 and efficient computation. \textit{Even small numerical errors} can result in
visible artifacts in rendered images.

Typical formulations involve scalar values such as $x$, indexed parameters
like $I_{i_j}$ , and exponential terms $e^{x^2}$
. These elements are combined to simulate
realistic lighting behavior.

\section{Core Components}
\subsection{Rendering Pipeline}
The rendering pipeline consists of the following stages:
\begin{itemize}
    \item Scene Definition
    \item Camera Configuration
    \item[-] Light Placement
\end{itemize}
\textbf{Important} Shading Model
\begin{itemize}
    \item[-]Diffuse Shading
    \item[-] Specular Highlights 
\end{itemize}
\begin{enumerate}
    \item Input Processing
    \item Ray Construction
         \begin{enumerate}
            \item How are you
            \item Secondary Rays
    \end{enumerate}
    \item Intersection Evaluation 
        \begin{itemize}
        \item Object Hit Test
        \end{itemize}
\end{enumerate}

\section{Mathematical Modeling}
Light interaction with surfaces can be expressed 
using the following equations.
\begin{equation}
L = I_{ij} \cdot \cos(\theta) + R^2 
\end{equation}

\[ k= \frac{ab}{cd} \]

\begin{align}
    f(x) &= x^2 + 2x + 1 \nonumber \\
    &= (x+1)^2
\end{align}




\end{document}
